% ======================================================================
% inst-txn-appendix.tex -- transactional-isolation instantiation (appendix).
%
% Carrier = schedules; commitments = ordering choices on the conflict graph
% (OCC validation order, MVTO timestamping, 2PL victim/lock-release); isolation
% levels (RC, SER, SI) are resolving subsets of ONE shared filtration. The
% depth theorem, the SER/SI support-depth incomparability, compositional
% insensitivity, SI worked depth, and protocol-specific structure all live
% here. Shares theorem/eqn labels with inst-txn-body.tex (only one of the two
% files is ever \input in a given build).
% ======================================================================
\section{Instantiation: Transactional Isolation}
\label{sec:transactions}

In a transactional system the non-monotone operation is \emph{conflict
resolution}: when concurrent transactions touch the same objects, the system
must decide which effects persist and in what apparent order. We instantiate
the determination model so that these decisions become explicit
\emph{commitments} over a conflict graph, leaving the operational interleaving
(reads, writes, arrivals) a free, work-free environment process. Two payoffs
follow. First, why/why-not under ambiguity becomes well-defined: a tuple can be
absent because of a serialization commitment, not merely a missing write.
Second, the classical isolation levels---read committed (RC), snapshot
isolation (SI), serializability (SER)---turn out to be \emph{different resolving
subsets of one shared filtration}, so their comparison reduces to a per-tuple
support-depth comparison (Proposition~\ref{prop:ser-si-incomparable}).

\subsection{The Carrier, Specification, and Commitment Basis}
\label{sec:txn-setup}

Fix a finite workload $W$: a set of transactions $T_1,\dots,T_n$ with their
read and write sets over a set of objects, together with a query transaction
$T_Q$. An \emph{outcome} is a decision trace: a per-transaction commit/abort
verdict together with an apparent serialization order on the committed
transactions, consistent with some execution of $W$. We take the output set
$O$ to be these decision traces and the output refinement order $\OrdO$ to be
\emph{sub-trace inclusion}: $o_1 \OrdO o_2$ when $o_2$ extends $o_1$ by
resolving the fate or relative order of further transactions without
contradicting any verdict already fixed. The acceptable relation
$S(W) = \Spec(W)$ admits exactly the traces that correspond to some
serializable (more generally, $L$-admissible, below) schedule of $W$; an
individual schedule is the application carrier.

The operations themselves---arrivals, reads, writes, and the interleaving the
scheduler happens to observe---are \emph{environment events}. They extend the
history and may \emph{introduce} conflict edges, but they do not by themselves
exclude any admissible outcome; in the vocabulary of
Definition~\ref{def:free} they are content-preserving free transitions and
carry no work and no depth. The fateful acts are the ordering decisions among
conflicting transactions.

\begin{definition}[Conflict and conflict graph]
  \label{def:txn-conflict}
  Two operations of distinct transactions $T_i \neq T_j$ \emph{conflict} if
  they access a common object and at least one is a write. Following
  Adya~\cite{adya1999weak} we label the induced edge
  $T_i \xrightarrow{\ell} T_j$ by its type
  $\ell \in \{\mathsf{ww}, \mathsf{wr}, \mathsf{rw}\}$: write--write
  ($\mathsf{ww}$), write--read ($\mathsf{wr}$), and read--write
  ($\mathsf{rw}$, an \emph{anti-dependency}). The \emph{conflict graph}
  $G(W)$ is the directed labeled multigraph on the active transactions carrying
  one such edge per conflict, oriented by operation precedence.
\end{definition}

\begin{definition}[Ordering commitment basis]
  \label{def:txn-basis}
  The commitment basis is the set of \emph{ordering commitments}
  \[
    \Basis = \{\,\varphi_{T_i \prec T_j} \mid T_i, T_j \text{ conflict in } W\,\}
      \;\cup\; \{\,\seal\,\},
  \]
  where $\varphi_{T_i \prec T_j}$ irrevocably records that $T_i$ is serialized
  before $T_j$ on their contested objects, and $\seal$ (Definition~\ref{def:free}'s
  fateful counterpart for a batch) closes a batch of pending requests so that no
  further transaction may join the current resolution round. One ordering
  commitment deletes from the residual every trace that orders the pair the
  other way (and every trace in which an induced forbidden cycle would commit
  both), hence it is fateful. A $\seal$ is fateful whenever the set of
  transactions it fixes forces a verdict that an unsealed state left open.
\end{definition}

\noindent
The abort of a transaction is \emph{not} a separate basis element: an abort is
\emph{entailed} by the ordering commitments once a forbidden cycle is detected,
exactly as a derived positive atom is entailed by a seal in the
$\mbox{Datalog}^\neg$ instantiation. All standard protocols realize $\Basis$:
OCC fixes order by \emph{validation order}, MVTO by \emph{timestamp
assignment}, and 2PL by the time of first lock release (or by deadlock
\emph{victim selection}); Section~\ref{sec:txn-protocols} shows these differ
only in which commitments the environment makes for the system.

\paragraph{Isolation levels as resolving subsets of one filtration.}
Because the basis and hence $\Dets$ are fixed by $W$, every isolation level
reads the \emph{same} layer-prefix filtration $\mathcal F$ of
Section~\ref{sec:algebra}; the levels differ only in which commitments are
\emph{forced}, i.e.\ which resolving subset of $\Dets$ each admits. Following
Adya~\cite{adya1999weak}:
\begin{itemize}[nosep]
  \item \emph{RC} imposes no constraint on cycles: all traces are admissible,
        so no ordering commitment is forced.
  \item \emph{SER} forbids every directed cycle in $G(W)$ regardless of edge
        labels: at least one transaction on each cycle must abort.
  \item \emph{SI} permits exactly those cycles containing two
        \emph{adjacent} $\mathsf{rw}$ anti-dependency edges, but enforces
        \emph{first-committer-wins} (FCW) on $\mathsf{ww}$ edges: of two
        concurrent writers of an object, only the first to request commit may
        commit.
\end{itemize}

\subsection{Determination Depth by Isolation Level}
\label{sec:txn-depth-section}

\begin{theorem}[Determination depth for transactions]
  \label{thm:txn-depth}
  Over the ordering basis $\Basis$, the support depth of transaction fate
  depends only on the protocol's resolution class:
  \begin{enumerate}[nosep,label=(\alph*)]
    \item \emph{No forbidden cycle (RC).} Every trace is admissible, so no
          commitment is forced and every outcome is $\OrdO$-comparable:
          support depth $0$.
    \item \emph{Fully reactive protocol.} Under any level $L$ that forbids some
          cycle but has \emph{no} scheduling discretion, each event is resolved
          deterministically by the current state and the arriving request; the
          determination has support depth $0$.
    \item \emph{Scheduling discretion.} Under such an $L$ \emph{with}
          discretion (batching, commit-order choice, or victim selection),
          whose forbidden cycle pattern recurs around a surviving
          discretionary transaction, the worst-case support depth is
          $\Theta(n)$ in the number of transactions. The \emph{per-batch} depth
          is exactly $2$: $\seal$ the batch, then $\seq$ a single
          (commuting) layer choosing a processing/validation order within it.
  \end{enumerate}
\end{theorem}
\begin{proof}
  \emph{(a)} If no cycle type is forbidden, validation always passes and no
  transaction need abort; all acceptable traces are comparable under $\OrdO$,
  so $Y_a \in \{\emptyset, \Dets\}$ for every fate observable and
  $\sdepth = 0$ by Definition~\ref{def:sdepth}.

  \emph{(b)} A fully reactive protocol processes each event by the current
  state alone. Witness: textbook OCC (validate at the commit request, abort iff
  a cycle exists) or textbook MVTO (assign the timestamp at arrival, force the
  abort entailed by it). The system exercises no choice, so $\Dets$ is a
  singleton up to free transitions and every support lies in
  $\mathcal F_0 = \{\emptyset,\Dets\}$.

  \emph{(c) Upper bound.} Each ordering commitment settles at least one
  transaction's fate; with $n$ transactions at most $n$ sequential commitments
  suffice, so no support exceeds depth $n$, giving $O(n)$.

  \emph{Lower bound (the $T_\infty$ witness).} Let $T_\infty$ be a long-running
  transaction that reads an object $x$ and writes an object $y$. Each fresh
  arrival $T_i$ writes $x$ and reads $y$, forming a two-edge $\mathsf{rw}$--$\mathsf{rw}$
  cycle with $T_\infty$:
  \[
    T_\infty \xrightarrow{\mathsf{rw}} T_i
    \qquad\text{and}\qquad
    T_i \xrightarrow{\mathsf{rw}} T_\infty .
  \]
  Under a level that forbids this cycle, each round forces a \emph{binary}
  decision: abort $T_\infty$ (resolving all future conflicts at once) or abort
  $T_i$ (letting $T_\infty$ survive to meet $T_{i+1}$). If the scheduler keeps
  $T_\infty$ alive for $n{-}1$ rounds, round $i$'s commitment depends on round
  $i{-}1$'s---had $T_\infty$ been aborted earlier, no later cycle would
  form---so the layers do not commute and the fate of $T_\infty$ is settled only
  by a depth-$(n{-}1)$ prefix. Hence $\sdepth = \Theta(n)$.

  \emph{Per-batch depth.} Within a sealed batch the only remaining discretion is
  the choice of a processing (validation / timestamp / victim) order $\pi$, a
  total order on the pending transactions; all aborts, commits, and read
  verdicts are \emph{entailed} by $\pi$. Thus a batch resolves in two
  non-commuting layers, $\seal \seq \pi$: the seal fixes \emph{which}
  transactions are in the round, and the order fixes \emph{which} survive. These
  do not commute (Remark~\ref{rem:layers}): the seal determines the vertex set
  on which $\pi$ then acts. Hence per-batch depth is $2$. Under 2PL the second
  layer is victim selection rather than order selection, but the seal is still
  required so that aborting a victim cannot admit a new request that forms a
  fresh cycle (Section~\ref{sec:txn-protocols}).
\end{proof}

\begin{example}[Per-batch depth under OCC with batch validation]
  \label{ex:txn-batch}
  Three transactions $T_1,T_2,T_3$ finish their read/write phases concurrently.
  The system seals the batch (layer~$1$): no further transaction joins this
  validation round. It then chooses a validation order $\pi$ (layer~$2$)---say
  $T_1 \prec T_2 \prec T_3$. Under $\pi$, $T_2$'s write set is checked against
  $T_1$'s; if they conflict, $T_2$ aborts---\emph{entailed} by $\pi$, not a
  separate commitment. Choosing $T_2 \prec T_1 \prec T_3$ instead may let $T_2$
  commit and abort $T_1$. The two layers do not commute, matching
  Theorem~\ref{thm:txn-depth}(c): depth~$2$ per batch.
\end{example}

\subsection{A Running Example}
\label{ex:txn-running}

We revisit the teaser of Example~\ref{ex:intro}. A relation $S(k,y)$ starts as
$\{S(0,b)\}$, and concurrent transactions
\[
  T_1\!:\textsf{ins } S(1,b),\quad
  T_2\!:\textsf{ins } S(2,d),\quad
  T_3\!:\textsf{del } S(2,d)
\]
run alongside the query transaction $T_Q\!: Q(y) \drule S(k,y)$. The relevant
conflicts are: $T_2 \xrightarrow{\mathsf{ww}} T_3$ (both touch $S(2,\cdot)$) and
the read--write edges relating $T_Q$ to $T_2,T_3$. Under SER the conflict graph
is acyclic, so all ordering commitments commute and resolution occupies a single
layer; two query-outcome classes of determinations arise, with representatives
\[
  D_{\mathsf{in}} = \varphi_{T_2 \prec T_Q} \cdot \varphi_{T_Q \prec T_3}
  \quad(\text{$T_Q$ observes $d$}),
  \qquad
  D_{\mathsf{out}} = \varphi_{T_2 \prec T_3} \cdot \varphi_{T_3 \prec T_Q}
  \quad(\text{$d$ absent}).
\]
Reading provenance through the support homomorphism of
Section~\ref{sec:support}, the why-support of $b$ is all of $\Dets$ while that
of $d$ is $\{D_{\mathsf{in}}\}$'s cell:
\[
  Y_b = \Dets, \qquad Y_d = \{\,D \in \Dets \mid D \text{ orders } T_2 \prec T_Q \prec T_3\,\}.
\]
Thus $b$ is \emph{robust} ($\sdepth(b) = 0$) and $d$ is \emph{contingent}
($\sdepth(d) = 1$): value $b$ appears under every ordering, whereas $d$ appears
exactly when $T_Q$ observes $T_2$'s insert before an effective delete.
Figure~\ref{fig:det-tables} summarizes the determination structure of this
example across isolation levels.

\subsection{Per-Tuple Isolation Sensitivity}
\label{sec:txn-sensitivity}

Vandevoort et al.~\cite{vandevoort2025mixed} certify \emph{individual
transactions} as isolation-insensitive. Provenance over the resulting database
enables a finer and a broader cut at once: a query may select a single output
tuple, or join tuples produced by several transactions. The filtration lifts
per-transaction certification to \emph{query results}.

\begin{proposition}[SER/SI support-depth incomparability]
  \label{prop:ser-si-incomparable}
  In the Adya conflict-graph model, the following two 2-transaction gadgets
  witness both directions of SER/SI support-depth incomparability for a tuple
  $t$ derived from a transaction $T_i$:
  \begin{enumerate}[nosep,label=(\alph*)]
    \item \emph{Write skew} ($\sdepth_{\mathrm{SER}}(t) > \sdepth_{\mathrm{SI}}(t)$).
          $T_A$ reads $x$ and writes $y$; $T_B$ reads $y$ and writes $x$, forming
          a two-edge $\mathsf{rw}$--$\mathsf{rw}$ cycle. SI permits the cycle:
          both commit in every SI determination, so $t$ is robust,
          $\sdepth_{\mathrm{SI}}(t) = 0$. SER forbids the cycle: one of $T_A,T_B$
          must abort, so $t$ is contingent, $\sdepth_{\mathrm{SER}}(t) > 0$.
    \item \emph{First-committer-wins forced abort}
          ($\sdepth_{\mathrm{SI}}(t) > \sdepth_{\mathrm{SER}}(t)$).
          $T_1$ writes $x$ and inserts $t$ into a relation $R$; $T_2$ writes $x$;
          there is no enclosing cycle. SER orders them on $x$ and commits both,
          so $t$ is robust, $\sdepth_{\mathrm{SER}}(t) = 0$. SI applies FCW on
          $x$ and forces one writer to abort, so $t$ holds only if $T_1$ wins:
          $\sdepth_{\mathrm{SI}}(t) > 0$.
  \end{enumerate}
\end{proposition}
\begin{proof}
  Both gadgets instantiate Definition~\ref{def:sdepth}. In (a) the SI resolving
  subset contains both-commit traces exclusively, so $Y_t = \Dets_{\mathrm{SI}}
  \in \mathcal F_0$ while under SER the surviving traces split by which
  transaction aborts, so $Y_t$ first appears at $\mathcal F_1$. In (b) the roles
  reverse: SER's acyclic graph fixes $t$'s presence in $\mathcal F_0$, whereas SI
  must choose an FCW winner, placing $Y_t$ at $\mathcal F_1$. Neither depth
  dominates the other across the two workloads, so the levels are incomparable.
\end{proof}

\begin{corollary}[Compositional isolation insensitivity]
  \label{cor:txn-insensitivity}
  If every transaction feeding a positive relational-algebra provenance query
  $Q$ is isolation-insensitive---equal support depth under SER and SI---then so
  is every output tuple of $Q$.
\end{corollary}
\begin{proof}
  Positive relational algebra computes output supports from input supports by
  the union and intersection of the support homomorphism
  (Proposition~\ref{prop:support-hom}); by
  Theorem~\ref{thm:filtration-respected} neither can raise support depth above
  the deepest input. If every input tuple has equal SER and SI support depth,
  the common bound transfers to each output tuple, so its support occupies the
  same filtration level under both levels.
\end{proof}

\begin{example}[Compositional insensitivity]
  \label{ex:txn-compose}
  Transactions $T_1,T_2,T_3$ insert into relations $R_1,R_2,R_3$; none
  participates in a write-skew or FCW pattern, so each is isolation-insensitive.
  The query $Q = R_1 \bowtie R_2 \bowtie R_3$ joins their outputs.
  Corollary~\ref{cor:txn-insensitivity} guarantees that every tuple of $Q$ is
  insensitive too, \emph{without} re-analyzing $Q$'s interaction with the
  conflict graph. This is the compositional extension of Vandevoort et al.'s
  per-transaction certification~\cite{vandevoort2025mixed}: the filtration
  certifies query results, not just transactions.
\end{example}

\begin{figure*}[t]
  \centering
  \footnotesize
  \setlength{\tabcolsep}{4pt}
  \renewcommand{\arraystretch}{1.15}
  \begin{tabular}{@{}lll@{}}
    \toprule
    \textbf{Isolation / case} & \textbf{Resolving determination} & \textbf{Support depth} \\
    \midrule
    Read committed &
    (none---no outcome contradicted) & $0$ \\
    SER, acyclic graph &
    $\{\varphi_{T_1 \prec T_2},\, \varphi_{T_2 \prec T_Q},\, \ldots\}$ (all commute) & $1$ \\
    SER, $d$ visible &
    $\varphi_{T_2 \prec T_Q} \cdot \varphi_{T_Q \prec T_3}$ & $1$ \\
    SER, $d$ absent &
    $\varphi_{T_2 \prec T_3} \cdot \varphi_{T_3 \prec T_Q}$ & $1$ \\
    SER, recurring cycle &
    $\varphi_{\mathsf{a}}(T_1) \seq \varphi_{\mathsf{a}}(T_2)$
    \;\emph{(Thm.~\ref{thm:txn-depth}c)} & $O(n)$ \\
    Snapshot isolation &
    $\seal \seq \varphi_{\mathrm{fcw}(S(5),T_6)}
      \seq \varphi_{\snap(T_Q)}$ & $O(n)$ \\
    \bottomrule
  \end{tabular}
  \smallskip
  {\footnotesize
    $\varphi_{T_i \prec T_j}$ = ordering;\;
    $\varphi_{\mathsf{a}}(T)$ = abort (entailed);\;
    $\{\cdots\}$ = commuting set;\;
    $\seq$ = non-commuting sequence.}
  \caption{Determination structure of the running example
    (Section~\ref{ex:txn-running}) across isolation levels, assuming scheduling
    discretion. Support depth is $0$ under RC or a fully reactive protocol, and
    $O(n)$ worst-case under SER or SI with discretion
    (Theorem~\ref{thm:txn-depth}).}
  \Description{Table showing resolving determinations and their support depth
    across isolation levels under scheduling discretion.}
  \label{fig:det-tables}
\end{figure*}

\subsection{Snapshot Isolation and the Fine-Grained SER/SI Separation}
\label{sec:txn-si}

Under a fully reactive SI implementation, snapshot assignment is deterministic
(each transaction reads the committed state at its start time) and FCW is
deterministic (whichever concurrent writer's commit request arrives first wins;
the other aborts). All verdicts are entailed by the environment, so support
depth is $0$. With discretion over commit ordering (which concurrent writer
commits first) or over batching (when snapshots are assigned), support depth
grows to $\Theta(n)$ by the same structure as SER.

\begin{proposition}[Worst-case SI support depth]
  \label{prop:txn-si-depth}
  Under snapshot isolation with commit-order discretion, the worst-case
  support depth is $\Theta(n)$.
\end{proposition}
\begin{proof}
  Mirror the $T_\infty$ witness of Theorem~\ref{thm:txn-depth}(c). Take $n$
  pairs $(T_i, T_i')$ in which $T_i$ and $T_i'$ both write object $x_i$ (an FCW
  conflict), and $T_{i+1}$ reads $x_i$ before writing $x_{i+1}$, so $T_{i+1}$'s
  snapshot depends on the FCW winner at $x_i$. At layer $i$ the system commits
  an FCW winner on $x_i$; this fixes which version $T_{i+1}$ sees, hence whether
  $T_{i+1}$ writes $x_{i+1}$, hence the FCW situation at layer $i{+}1$. The
  decisions do not commute, giving a depth-$n$ chain; the $O(n)$ upper bound is
  as before. Per batch the depth is again $2$ ($\seal \seq$ commit-order
  selection).
\end{proof}

\paragraph{Worked SI depth.}
Suppose the database holds $S(5,e)$, and $T_6\!:\textsf{write } S(5,e')$,
$T_7\!:\textsf{write } S(5,e'')$ run concurrently; a query transaction $T_Q$
reads key $5$ after both request commit. Under SI only one of $T_6,T_7$ may
commit (FCW on key $5$); the other aborts. With commit-order discretion this is
a genuine commitment: choosing $T_6$ makes $T_Q$ read $e'$, choosing $T_7$ makes
it read $e''$. The determination has depth $2$,
\[
  D_{\mathrm{SI}} = \seal \seq \varphi_{\mathrm{fcw}(S(5),\,T_6)},
\]
and under it $T_Q$ reads $S(5,e')$ with conditioned provenance
$P_{D_{\mathrm{SI}}}(e') = x_6$; the alternative
$\varphi_{\mathrm{fcw}(S(5),T_7)}$ yields the contingent $e''$. Here $e'$ has
$\sdepth = 1$ (its why-support is the FCW-winner cell at layer $1$ after the
seal), illustrating Proposition~\ref{prop:txn-si-depth} at $n = 1$.

\paragraph{Characterization of the SER/SI gap.}
Fix a workload $W$ with conflict graph $G$. For each level
$L \in \{\mathrm{SER},\mathrm{SI}\}$ the determination space $\Dets$ and
filtration $\mathcal F$ are built over the same $\Basis$ but with $L$'s
admissibility constraint selecting the resolving subset. A tuple $t$ derived
from $T_i$ has $\sdepth_L(t) = 0$ iff $T_i$ commits in every resolving
determination under $L$ (robust) and $\sdepth_L(t) > 0$ iff $T_i$'s fate varies
(contingent). The two agree exactly when $T_i$ lies on no cycle with two
adjacent $\mathsf{rw}$ anti-dependency edges (so SER and SI agree on whether a
cycle must break; Adya~\cite{adya1999weak}, Fekete et
al.~\cite{fekete2005making}) \emph{and} in no $\mathsf{ww}$ conflict without an
enclosing cycle (so they agree on whether both writers commit). These are
precisely the robustly portable transactions
of~\cite{vandevoort2025mixed}; outside this class, the two directions of
Proposition~\ref{prop:ser-si-incomparable} apply.

\paragraph{Certificates and semantic-change queries.}
A \emph{certificate} over-approximates a commitment's effect: it records which
outcomes the commitment excludes, admitting possibly more but never fewer, and
is exactly a why/why-not certificate in the sense of
Section~\ref{sec:filtration}. Under SER a certificate need only record the
ordering predicates sufficient to justify the observation; under SI, answering
semantic-change queries---e.g.\ whether an SI outcome admits a serial
explanation---requires retaining the snapshot structure as well. Certificates
are thus query-parametric: richer questions demand deeper prefixes, in line with
support depth.

\subsection{Protocol-Specific Determination Structures}
\label{sec:txn-protocols}

Theorem~\ref{thm:txn-depth} gives protocol-agnostic bounds via the $T_\infty$
witness. The three standard protocols instantiate $\Basis$ differently but share
the $\seal + \text{resolution-layer}$ structure and the same depth bounds; their
bases are pairwise incomparable in \emph{which} act the system commits.

\paragraph{OCC (validation-order basis).}
A sealed batch of pending commit requests is resolved by a chosen validation
order; transactions validated earlier commit (if no cycle among those already
committed), later ones may find a cycle and abort. The abort is entailed by the
order, matching Example~\ref{ex:txn-batch}. Per-batch depth $2$: $\seal \seq$
validation order.

\paragraph{2PL (victim-selection basis).}
Lock acquisition and release are fully reactive (entailed by the arrival order
of requests); the only discretion is \emph{victim selection} when a deadlock
cycle forms. Without sealing, victim choices cascade---aborting a victim
releases locks, admitting new requests that may form new cycles, giving
$\Theta(n)$ depth. Sealing the batch (refusing new lock requests) makes all
deadlock cycles in the wait-for graph present simultaneously, so victim choices
across distinct cycles commute. Each lock acquisition that blocks a conflicting
request realizes an ordering commitment $\varphi_{T_i \prec T_j}$; a deadlock
cycle means the acquisitions already made are jointly inconsistent with all
committing, and the chosen victim is an entailed abort. Per-batch depth $2$:
$\seal \seq$ one commuting layer of victim selections. Fully reactive (immediate
deadlock resolution with the requestor as victim) gives depth $0$;
victim-selection discretion gives $O(n)$ (Theorem~\ref{thm:txn-depth}).

\paragraph{MVTO (action-ordering basis).}
Under multiversion timestamp ordering~\cite{reed1978naming}, each transaction
receives a timestamp at start; reads return the latest version with timestamp
$\le$ the reader's, and a write is rejected (its transaction aborted) if a
higher-timestamped transaction has already read an earlier version. In the fully
reactive case the primary commitment is timestamp assignment: assigning
$ts_i < ts_j$ fixes $T_i \prec T_j$ on every shared object at once---a single
ordering commitment with wide reach. With batching, the system instead chooses a
processing order over individual \emph{actions} (starts, reads, writes), which
subsumes timestamp assignment and determines which reads see which versions and
hence which transactions abort. MVTO never blocks, so it is deadlock-free; the
determination again factors as $\seal \seq$ an action-ordering layer, at action
rather than transaction granularity. Per-batch depth $2$.

\paragraph{Shared structure, incomparable bases.}
All three realize $\Basis$ and reach the same $\Theta(n)$ worst case and
per-batch depth $2$, but their bases are incomparable: OCC commits a
transaction-ordering act at validation time, 2PL commits victim choices
distributed throughout execution, and MVTO commits an action-ordering act via
timestamps. The common $\seal + \text{resolution-layer}$ shape is what lets the
support-depth analysis---and Theorem~\ref{thm:txn-depth}---apply uniformly,
independent of protocol.
